% SPDX-FileCopyrightText: © 2026 Idiap Research Institute <contact@idiap.ch>
%
% SPDX-FileContributor: Andre Freitas <andre.freitas@idiap.ch>
% Neuro-symbolic AI Group
%
% SPDX-FileContributor: Danilo Gusicuma <danilo.gusicuma@idiap.ch>
% Neuro-symbolic AI Group
%
% SPDX-License-Identifier: MIT
% Generated from step-explanations.json
\paragraph{Step-by-step reconstruction.}
Prove Euclidean-norm decay for the symmetric combinatorial Laplacian of a finite connected undirected graph on mean-zero initial data, assuming its first positive eigenvalue is at least c greater than zero.
The following is a source-guided illustrative reconstruction. Internal procedures are prescriptions; the native declarations do not establish their execution.
\begin{enumerate}[leftmargin=*,label=\arabic*.]
\item \textbf{Construct the exact conditional theorem} (\texttt{d0}).
\begin{description}[style=standard,labelindent=0pt,leftmargin=5.5em,labelwidth=5em,labelsep=0.5em,align=left,font=\normalfont\bfseries,itemsep=0.15em,parsep=0pt,topsep=0.3em]
\item[Purpose] Derive the exact queued obligation.
\item[Input] item.premises: queued exact\_\allowbreak{}statement; connected-graph, subspace and gap premises
\item[Method] Order lemmas without circular dependence and propagate constants, norms and quantifiers through each inference.
\item[Check] Every inferential step must identify its premises and preserve the theorem's scope.
\item[Output] Exact claim; retained premises and undischarged formal obligations
\item[Next] Refutation receives this claim; a missing lemma remains visible rather than silently closed.
\item[Example] For du/dt=-Lu on the connected graph, expand the mean-zero vector in positive-eigenvalue modes. Each coefficient evolves by exp(-lambda t), giving Euclidean-norm decay at rate at least c.
\end{description}
\item \textbf{Search for a defeater} (\texttt{r0}).
\begin{description}[style=standard,labelindent=0pt,leftmargin=5.5em,labelwidth=5em,labelsep=0.5em,align=left,font=\normalfont\bfseries,itemsep=0.15em,parsep=0pt,topsep=0.3em]
\item[Purpose] Challenge edge cases without weakening the claim.
\item[Input] Exact claim; gap premise; open obligations
\item[Method] Test small, degenerate and adversarial cases; classify a failed step as a false statement, wrong quantifier, model mismatch or missing technique.
\item[Check] Favorable examples cannot substitute for a missing spectral argument.
\item[Output] Refutation assessment: scoped surviving theorem or localized defect
\item[Next] Calibration retains this finding with the exact claim; repair must name any changed premise.
\item[Example] The constant mode violates unrestricted decay but is excluded by the mean-zero premise. Without the gap premise, cycle families defeat a common positive rate.
\end{description}
\item \textbf{Calibrate the proof disposition} (\texttt{e0}).
\begin{description}[style=standard,labelindent=0pt,leftmargin=5.5em,labelwidth=5em,labelsep=0.5em,align=left,font=\normalfont\bfseries,itemsep=0.15em,parsep=0pt,topsep=0.3em]
\item[Purpose] Distinguish proof discharge from conditional or blocked work.
\item[Input] Exact theorem claim + its refutation assessment
\item[Method] Check the lemma chain and retained obligations; issue a bounded disposition and failure certificate when a bridge remains missing.
\item[Check] A proof gate cannot discharge an obligation merely because refutation found no counterexample.
\item[Output] Claim; finding; proof gate; optional failure certificate
\item[Next] Original policy routes accepted work or failure; the native return is no terminal proof acceptance act.
\item[Example] Accept only the explicitly justified conditional spectral estimate; if spectral decomposition has not been supplied, retain that proof debt in the disposition.
\end{description}
\end{enumerate}
\paragraph{Changed-input contrast.}
If the queue requests all initial data, a nonzero constant vector defeats positive-rate decay. The theorem must retain the mean-zero restriction; an assessment of that restricted theorem cannot certify the broader one.
